% !TeX root = ../main.tex
\section{Discussion}
\label{sec:discussion}

\textbf{Choosing an operating point.}
No single metric captures deployability. \emph{jitter\_high} formally lies on the
bandwidth--accuracy Pareto frontier but imposes 221\,ms mean latency---unacceptable
for interactive QUIC. MTU padding maximizes privacy but roughly quadruples bytes
sent---infeasible for backbone operators. \emph{jitter\_low} occupies a sweet spot:
statistically distinguishable from baseline yet only 1.7 macro-F1 points lower, with
11\,ms latency and no bandwidth penalty. We recommend it as the default proxy policy
when bandwidth is constrained; stronger jitter tiers suit delay-tolerant bulk flows.

\textbf{Transformer vs.\ BiLSTM.}
Global attention explains higher robustness to timing perturbation: distant packets
remain correlated even when local IPT order is noised. BiLSTM's sequential inductive
bias amplifies jitter effects (\emph{jitter\_high} gap $\approx$18.7 percentage
points accuracy). For padding-only defenses the ranking reverses---BiLSTM slightly
wins on \emph{linear128}---so defender--attacker co-design matters. Reporting a single
``best architecture'' without specifying defense is misleading.

\textbf{Limitations.}
(i) Attackers are not retrained on obfuscated data---an adaptive adversary would
partially recover accuracy. (ii) Single-network, single-month evaluation; seasonal
drift and cross-ASN generalization are unverified. (iii) Deterministic obfuscation is
predictable; adversaries who model the proxy transform may craft robust features.
(iv) Class labels in per-class reports are numeric IDs; mapping to application names
requires the released label encoder.

\textbf{Future work.}
Adaptive attacker retraining, composable stochastic defenses, integration with QUIC
PADDING frames and DATagram coalescing, and live testbed validation on a QUIC proxy
(MASQUE/Oblivious HTTP) are natural extensions.

%==============================================================================
